\subsection{素域}

有理整数环 Z 的分式域称为有理数域，记为 Q（I，第 117 页）。对每个素数 p，商环 \(\mathbf{Z}/(p)\) 是一个有限的 \(^{1}\) 、含 p 个元素的域，下文记为 \(F_{p}\) 。域 Q 是无限的，因为它包含 Z，因此不同构于任何域 \(F_{p}\) 。若 p 与 \(p'\) 是不同的素数，则域 \(F_{p}\) 与 \(F_{p'}\) 具有不同的基数，因此不同构。

定义 1. --- 一个域称为素域，如果它同构于 Q 或同构于某个域 \(F_{p}\) .

Q 的每个子域都包含环 Z，从而包含 Z 的分式域 Q；\(F_{p}\) 的每个子环都必然等于 \(F_{p}\) 。因此，素域的每个子域必然等于它本身（参见下面定理 1 的推论 2）。设 P 是一个素域，A 是一个环；若 f 与 \(f'\) 是 P 到 A 中的两个同态，则 \(x \in P\) （满足 \(f(x) = f'(x)\)）组成的集合是 P 的一个子域，因此根据上述结论必有 \(f = f'\) 。特别地，素域的唯一的自同态是恒等映射。

定理 1. --- 设 A 是一个环；假设 A 存在一个子域。则 A 有唯一的子域 P，它是一个素域。此外，P 包含在 A 的中心以及 A 的每个子域中。

设 K 为 A 的子域，C 为 A 的中心，并令 \(K' = K \cap C\) ；则 \(K'\) 是 A 的子域。设 f 为 Z 到 A 的唯一同态，p 为其核。A 的每个子环，特别是 \(K'\) ，包含 \(f(\mathbf{Z})\) ；因此理想 p 是素理想（I，第 116-117 页）。如果 \(\mathfrak{p} = (0)\) ，从 Z 到 \(K'\) 的同态 f 是单射；因此它（I，第 116 页）扩展为一个同构 \(\bar{f}\)（把 Q 映到 \(K'\) 的一个子域 P 上）。若 \(\mathfrak{p} \neq (0)\) ，存在一个严格正整数 p，使得 \(\mathfrak{p} = (p)\) （I，第 111 页）；如果我们有 p = ab，其中 a \textgreater{} 1 且 b \textgreater{} 1，这将意味着 \(a \notin p\) , \(b \notin p\) 且 \(ab \in p\) ，与 p 为素数的事实相矛盾。因此数 p 是素数，并且通过取商，f 定义了 \(F_p = Z/p\) 到 \(K'\) 的一个子域 P 上的一个同构。在这两种情形下，P 都是 A 的一个子域，且包含在 A 的中心 C 中，并且它是一个素域。设 L 是 A 的一个子域；则 \(P \cap L\) 是 P 的一个子域，由于 P 是素域，我们有 \(P \cap L = P\) ，从而 \(P \subset L\) 。若 \(P'\) 是 A 的一个子域且是素域，那么根据前面所述， \(P \subset P'\) ，从而 \(P = P'\) ，因为 \(P'\) 是素域。

推论 1. --- 设 K 是一个域。K 中存在唯一的素域子域，并且它是 K 的最小（在包含关系下）子域。

推论 2. --- 一个域是素域，当且仅当它不包含除自身以外的任何子域。

